\documentclass[oneside,a4paper,12pt]{article}

\usepackage[english,brazilian]{babel}
\usepackage[utf8]{inputenc}
\usepackage[T1]{fontenc}

\usepackage[top=20mm, bottom=20mm, left=20mm, right=20mm]{geometry}
\usepackage{framed}
\usepackage{booktabs}	   		% Pacote para deixar tabelas mais bonitas.
\usepackage{color}				% Pacote de Cores
\usepackage{hyperref}			% Pacotes para Hiperlinks
\usepackage{graphicx}			% Pacote de imagens

\definecolor{shadecolor}{rgb}{0.8,0.8,0.8}

\newcommand{\universidade}{Universidade Federal do Rio Grande do Norte (UFRN) \par Centro de Ciências Humanas, Letras e Artes (CCHLA)}
\newcommand{\curso}{Geografia}
\newcommand{\professores}{Nome de Aluno 1; Nome de Aluno 2; Nome de Aluno 3; Nome de Aluno 4; Nome de Aluno 5; Nome de Aluno 6}
\newcommand{\disciplina}{Metodologia da Ciência}
\newcommand{\tema}{A história da filosofia como desenvolvimento da ideia}
\newcommand{\turma}{1º Periodo}
\newcommand{\data}{28/05/2014}
\newcommand{\tempodeaula}{30 minutos}


\begin{document}

  \pagestyle{empty}

	\begin{center}
	
	  \universidade
       \par
	  Curso de \curso
	  \par
	  \vspace{10pt}
	  \LARGE \textbf{Plano de Seminário}

	\end{center}

  \vspace{10pt}

	\begin{tabular}{ |l|p{12cm}| }

	  \hline
	  \multicolumn{2}{|l|}{\textbf{Dados de Identificação}} \\
	  \hline
	  Professores:         &    \professores           \\
	  \hline
	  Disciplina:        &    \disciplina          \\
	  \hline
	  Tema:              &    \tema                \\
	  \hline
	  Turma:             &    \turma               \\
	  \hline
	  Data:              &    \data                \\
	  \hline
	  Duração da aula:   &    \tempodeaula         \\
	  \hline

	\end{tabular}

  \begin{snugshade}
  \section{Objetivos} % a serem alcançados pelos alunos e não pelo professor. Podem ser divididos em gerais e específicos. 
  \end{snugshade}
  
    \subsection{Geral} % projeta resultado geral relativo a execução de conteúdos e procedimentos.

	Entender como Hegel via a relação entre o desenvolvimento histórico e o desenvolvimento do espírito humano, além de conhecer um pouco sobre a influência de Hegel na política mundial e na filosofia.
    
    \subsection{Específicos} % especificam resultados esperados observáveis (geralmente de 3 a 4).
    
      \begin{itemize}
      
		\item Conhecer algumas definições relativas aos principais termos usados;
		\item Entender a ideia de processo histórico e as diferenças relativas a algumas culturas observadas por Hegel;
		\item Entender a racionalidade do processo histórico e o papel da consciência e da liberdade em seu desdobramento;
		\item Conhecer a repercussão do pensamento de Hegel na política mundial e alguns desdobramentos contemporâneos.
      
      \end{itemize}
  
  \begin{snugshade}
  \section{Conteúdos} % conteúdos programados para a aula organizados em tópicos (de 4 a 8).
  \end{snugshade}

      \begin{itemize}
      
		\item A Filosofia da História de Hegel, numa casca de nós; %---> inclui definição dos termos 
		\item O início do processo histórico;
		\item A racionalidade e a história da busca pela liberdade;
		\item O legado de Hegel.
      
      \end{itemize}

  \begin{snugshade}
  \section{Procedimentos metodológicos} % estratégias relevantes adotadas para alcançar os objetivos.
  \end{snugshade}

	Apresentação expositiva e dialogada. Dinâmicas de Ensino-aprendizagem. Problematização de Estudos de Casos.
	  
  \begin{snugshade}
  \section{Recursos didáticos} % quadro, giz, retro-projetor, filme, música, quadrinhos, etc.
  \end{snugshade}
  
    \begin{itemize}

	  \item Computador, projetor multimídia, pincel e quadro.
	  
    \end{itemize}
  
  \begin{snugshade}
  \section{Avaliação} % pode ser realizada com diferentes propósitos (diagnóstica, formativa e somativa). Interessante explicitar a atividade avaliativa e os critérios de correção.
  \end{snugshade}
  
  Os alunos deverão demonstrar uma compreensão suficiente para, por exemplo, responder questões como:

    \begin{enumerate}
	  
	  \item Qual a relação entre espírito, mente e auto-consciência?
	  \item Qual a importância da ideia de um processo histórico capaz de ser compreendido pela razão?
	  \item Porque é importante que o espírito tenha consciência de si e o que isso implica em relação à liberdade e ao processo histórico?
	  	  
	\end{enumerate}

%\cleardoublepage

% Referências bibliográficas


\begin{thebibliography}{}

\bibitem{French2009}
FRENCH, Steven
\newblock \textbf{Ciência}
\newblock Conceitos-chave em Filosofia. Trad.: André Klaudat. Porto Alegre: Artmed, 2009.

\bibitem{Shinn2008}
SHINN, Terry. 
\newblock Regimes de produção e difusão de ciência: rumo a uma organização transversal do conhecimento.\newblock \textbf{Scientiae Studia}
\newblock v. 6, n. 1, p. 11-42, 2008. 
\newblock Disponível em: <\url{http://www.scielo.br/scielo.php?script=sci_pdf&pid=S1678-31662008000100002&lng=pt&nrm=iso&tlng=pt}>. Acesso em: 16 mai. 2014.


\end{thebibliography}

\pagebreak

  \begin{snugshade}
  \section*{Anexo A - Mapa Conceitual} % Mapa conceitual do tema
  \end{snugshade}


\begin{center}
	\includegraphics[width=18cm]{mapa.png} 
\end{center}


\end{document}